% ---------------------------------------------------------------------------
% Section 3 -- Problem Setting and Threat Model
%
% Job: make the setting precise enough that Section 4 can be formal.
%   (a) the MCP mechanics,          (b) actors and who is protected from whom,
%   (c) how harm arises,            (d) the server's observation set,
%   (e) assumptions and scope,      (f) the decision problem, stated formally.
%
% Budget: ~650 words + one small table.
%
% TERMINOLOGY: "resource" is reserved for r, the thing a call touches. The
% server is the "protected party", never "the protected resource", and the word
% "asset" does not appear here. Sec.4/5 still say "asset register" and
% "(tool, asset) cell" -- one sweep would make the paper uniform.
%
% THE ROUTES-TO-HARM LIST is now ONE sentence with grouped citations: Sec.1 P2
% and Sec.2 P1/P2 already carry the same six keys, and this was the third
% telling of the same list.
%
% Sec.4.7 (sec:enforcement) currently re-states the risk-adaptive AC pattern
% with the same three keys. This section now establishes it, so that sentence
% in Sec.4.7 can be reduced to a back-reference.
% ---------------------------------------------------------------------------
\section{Problem Setting and Threat Model}
\label{sec:threat}

\textbf{The setting.} An MCP server publishes a catalog of tools over a uniform
JSON-RPC channel. A client issues \texttt{tools/list} to obtain it --- each
tool's name, natural-language description, and JSON argument schema --- then
invokes one with concrete
arguments~\cite{anthropic2024mcp,mcp2026spec}. The tools are an interface, not
the resource: behind them sit the organization's own repositories, calendars,
message histories and payment rails, and the server, not the agent, is the
process that touches them. Every invocation is an action the server performs on
its own resources at an agent's request.

\textbf{Actors.} The \emph{server} is the protected party: it absorbs the
consequence of anything done to those resources. The \emph{agent}
is the risk source, an LLM-driven client choosing which tool to call and with
which arguments. The \emph{organization} deploys the server, owns the
resources, and publishes the policy that states what losing them would mean. The
calling population is open where it counts: a deployment authorises a
connection, not an implementation, so the server can inspect neither the model
nor the plan behind a request.

\textbf{How harm arises.} Harm arises even though the agent is authenticated and
authorised for the tool it calls --- through hallucinated calls, risk the agent
can articulate but does not act on, capability exceeding what the task
needs, injected content, and objectives that diverge from the
principal's~\cite{yin2025reasoningtrap,tang2025riskknowledge,li2025privilege,%
huang2026auditing,zhan2024injecagent,lynch2025misalignment}. These differ in
cause but not on the wire: each produces a well-formed, schema-conforming,
authenticated call from a permitted agent to a permitted tool, and no
protocol-level check separates them. What varies is what executing it would do.
The opposite direction, a malicious server attacking the agent, is out of scope.

\textbf{What the server observes.} At decision time --- after a call arrives,
before it executes --- the server sees the caller's identity, already consumed by
authorization, and the five quantities in Table~\ref{tab:observables} that
describe the invocation itself. Identity is absent from the five because it does
not separate the cases: a correct agent, a mistaken one and a manipulated one
present the same credentials over the same
interface~\cite{huang2026caller,mellafe2026capability}. Nor does the server
observe \emph{intent}: it cannot see the user's instruction, and cannot tell
whether an argument naming, say, a \$1{,}000{,}000 transfer is a typo, an
injection, or a sanctioned operation. The score below is therefore a function of
those five alone, never of what the agent meant.

\begin{table}[t]
  \caption{What an MCP server can observe about an invocation before executing
  it. Intent is not among them.}
  \label{tab:observables}
  \begin{tabular}{@{}cl@{}}
    \toprule
    & Observable \\
    \midrule
    $t$ & the tool named in the request \\
    $r$ & the resource(s) that request resolves to \\
    $a$ & the argument values, valid against the tool's schema \\
    $x$ & session context: the calls already made on this connection \\
    $h$ & the calling agent's own history of prior invocations \\
    \bottomrule
  \end{tabular}
\end{table}

\textbf{Assumptions.} \emph{(i)} Authentication and transport security are
sound and orthogonal: identity is given, and we ask only what that caller may
safely do. \emph{(ii)} The organization publishes a written
classification policy of the kind security categorization standards
define~\cite{nist2004fips199,nist2008sp80060} --- named classes defined by
adverse impact, plus rules for recognising a resource's class --- and registers
the resources its server touches while assigning a sensitivity number to none of
them. We take this to be the realistic case; rules, unlike an inventory, still
classify resources created after the policy was written. \emph{(iii)} A call
resolves to rows of that register by its tool and, where several are reachable,
by its arguments. \emph{(iv)} The server can refuse or hold a call, so
enforcement is in-band rather than advisory.

\textbf{The decision problem.} Write an invocation as the whole tuple
$c = \langle t, r, a, x, h \rangle$ --- not the tool alone. Access control
supplies a predicate $\mathrm{perm}(c) \in \{0,1\}$ and already refuses what it
rejects, so the population of interest is exactly the calls with
$\mathrm{perm}(c) = 1$. The problem is to compute a graded
consequence score $\rho(c)$ that orders them by how much damage executing them
could do, and to gate on it against thresholds $\theta_{r} \le \theta_{d}$ the
deploying organization sets --- conditioning authorization on a computed risk
value, as risk-adaptive access
control does~\cite{cheng2007fuzzymls,kandala2011radac,atlam2020riskbased}:
\begin{equation}
  d(c) =
  \begin{cases}
    \textsf{allow}  & \rho(c) < \theta_{r} \\
    \textsf{review} & \theta_{r} \le \rho(c) < \theta_{d} \\
    \textsf{deny}   & \rho(c) \ge \theta_{d}.
  \end{cases}
  \label{eq:decision}
\end{equation}
With $\rho$ reported as a discrete band the thresholds sit on adjacent bands and
\textsf{review} is the band at $\theta_{r}$; \S\ref{sec:eval} operates a single
flag threshold, swept across bands. Two properties make $\rho$, not
$d$, the object to design. It must be ordinal rather than binary, so an operator
sets where escalation falls instead of inheriting the designer's cut
point~\cite{cao2024mad}; and it must be auditable, decomposed into named factors
a security team can argue with, since $\theta_{r}$ and $\theta_{d}$ are an
organizational judgement, not a tuned hyperparameter.
Section~\ref{sec:framework} constructs $\rho$.
